\documentclass[11pt]{article}
\usepackage[margin=1in]{geometry}
\usepackage{amsmath,amssymb,amsthm,mathtools,bm}
\usepackage{microtype}
\usepackage[T1]{fontenc}
\usepackage{lmodern}
\usepackage{booktabs,array,enumitem,caption}
\usepackage[hidelinks]{hyperref}
\setlist{nosep}
\newtheorem{theorem}{Theorem}[section]
\newtheorem{corollary}[theorem]{Corollary}
\newtheorem{lemma}[theorem]{Lemma}
\newtheorem{proposition}[theorem]{Proposition}
\theoremstyle{definition}
\newtheorem{definition}[theorem]{Definition}
\newtheorem{assumption}[theorem]{Assumption}
\newtheorem{model}[theorem]{Modeling Assumption}
\theoremstyle{remark}
\newtheorem{remark}[theorem]{Remark}

\newcommand{\dd}{\mathrm{d}}
\newcommand{\R}{\mathbb{R}}

\title{\textbf{Hysterical Conformal Cyclic Cosmology}\\[0.35em]
\large Infrared electromagnetic criticality and an active big-bang crossover}
\author{Andrew Bruce Boyd}
\date{Working mathematical formulation --- September 2026}

\begin{document}
\maketitle
\noindent\textit{Computational note.}
Proofs, exploratory experiments, and paper writing were assisted by generative AI.
Mathematical theorems stated in this paper are formally verified in Lean~4 in the
companion specifications, as faithful ports of the TeX arguments at the abstractions
recorded there, except results explicitly tagged as \emph{literature hypotheses}
(standard theorems cited from the literature and not re-proved here) or
\emph{physical inputs} (modeling assumptions outside mathematics).

\begin{abstract}
The late-universe gravitational response of this series forces dilution-driven criticality
and a vacuum-like attractor.
Here we isolate a companion late-time channel: infrared electromagnetic response in an
asymptotically de~Sitter background.
A minimal helical three-mode linearisation has an exact characteristic polynomial; under
positive parameters the stationary instability is infrared-first and occupies a finite band
of physical wavenumber.
A gauge-invariant scalar--Maxwell kinetic completion then supplies a pump for finite-$k$
photons.
Reciprocal conformal matching across a crossover surface maps outgoing de~Sitter scale-factor
asymptotics to an incoming radiation-era law $a\propto t^{1/2}$, while regular conformal
Maxwell data appear cold on the outgoing side and hot on the incoming side.
We call the resulting picture \emph{Hysterical Conformal Cyclic Cosmology} (HCCC):
geometrically akin to Penrose's conformal cyclic cosmology, but driven by an
\emph{active} electromagnetic phase transition rather than a passive conformal
identification.
That crossover is the big-bang-like event of the next aeon.
We separate algebraic theorems from modelling closures and from open dynamical questions
(saturation amplitude, junction regularity, thermalisation).
\end{abstract}
\tableofcontents
\newpage

\section{Scope and status}
\label{sec:scope}
The response-theory paper isolates the Hysterical Response as a residual of an
already-existing interaction under exponentially stable reduced dynamics~\cite{boyd2026response}.
The late-universe continuation shows that homogeneous dilution of a conserved density,
together with a relaxation gap vanishing as a positive power of that density, forces a
unique criticality scale and a positive vacuum-like attractor~\cite{boyd2026late}.
That same late-universe paper records a sharp electromagnetic caveat: exact charge
conservation forbids generating a homogeneous net charge from $Q=0$, while finite-wavenumber
electromagnetic segregation remains allowed.

The present paper develops that electromagnetic window into a cosmological mechanism.
In deep late-time de~Sitter, an infrared electromagnetic response channel can become
supercritical.
The resulting phase transition pumps conformal radiation data and, under a reciprocal
conformal matching hypothesis of Penrose--Tod type~\cite{Penrose2006,Penrose2010,Tod2003},
supplies the initial radiation content of a subsequent aeon.
The claim is interpretive but sharp:
\begin{equation}
\label{eq:claim}
\boxed{\begin{gathered}
\text{IR EM Hysterical criticality}
\;+\;
\text{photon pump}
\;+\;
\text{reciprocal conformal matching}\\
\Longrightarrow\quad
\text{active big-bang-like crossover to a new aeon.}
\end{gathered}}
\end{equation}
We call this package Hysterical Conformal Cyclic Cosmology (HCCC).
It is deliberately compared with Penrose's conformal cyclic cosmology
(CCC)~\cite{Penrose2006,Penrose2010,Penrose2013CCC}: both identify the remote future of one
aeon with the Big Bang of the next through conformal geometry.
The distinguishing thesis is dynamical.
Standard CCC is \emph{passive}: after massive degrees of freedom become irrelevant, conformal
invariance licenses a smooth identification across the crossover~\cite{Penrose2010,Tod2003}.
HCCC is \emph{active}: a Hysterical infrared electromagnetic instability drives a phase
transition that produces the conformal radiation data of the next aeon.

Four logical layers must be kept separate.
\begin{center}
\begin{tabular}{@{}>{\bfseries}p{0.24\linewidth}p{0.68\linewidth}@{}}
\toprule
Theorem & Algebraic or elementary-analytic consequence of the stated hypotheses.\\
Modelling assumption & Scalar--Maxwell kinetic completion; reciprocal conformal ansatz; FLRW asymptotics.\\
Physical input / literature & Expanding late-time de~Sitter; 4D Maxwell conformal invariance; CCC geometry.\\
Not claimed & Derivation of the saturation amplitude; full Einstein--Maxwell--response junction;
entropy production; observational CCC rings/Hawking points; proof that Nature realises HCCC.\\
\bottomrule
\end{tabular}
\end{center}

\section{Inherited late-universe setting}
\label{sec:inherited}
We work on a spatially flat FLRW background with scale factor $a(t)$ and Hubble rate
$H=\dot a/a>0$.
In the late-universe gravitational channel of~\cite{boyd2026late}, dilution forces a unique
criticality scale $a_c$ under a density-vanishing gap law, after which a pitchfork attractor
yields a positive asymptotic density with $w\to-1$ from below.
That package is inherited here only as background motivation: once a vacuum-like term
dominates, the outgoing geometry is asymptotically de~Sitter,
\begin{equation}
\label{eq:dS}
a_-(\eta)=\frac{1}{H_-|\eta|}+o(|\eta|^{-1}),
\qquad \eta\to 0^-,
\end{equation}
in conformal time $\eta$ with constant $H_->0$.
We do not re-prove eventual gravitational criticality.

The electromagnetic distinction of~\cite{boyd2026late} remains in force.
Exact charge conservation still forbids a homogeneous net-charge instability from $Q=0$.
The mechanism below is therefore a field / finite-wavenumber electromagnetic response, not
creation of cosmic net charge.

\section{Linear infrared electromagnetic instability}
\label{sec:linear}
For one helical Fourier mode with physical wavenumber $q=k/a$ and helicity
$s\in\{-1,+1\}$, consider the three-variable linear response model
\begin{align}
\label{eq:linEM}
\dot E &= -2HE + s q B + \alpha h,\\
\dot B &= -2HB - s q E,\\
\dot h &= \sigma E - \gamma h.
\end{align}
Here $H>0$ is treated as locally constant (de~Sitter freezing), $\gamma>0$ is the response
relaxation rate, and only the product $A:=\alpha\sigma$ enters the spectrum.
The Jacobian is
\begin{equation}
\label{eq:J}
J=\begin{pmatrix}-2H&sq&\alpha\\-sq&-2H&0\\\sigma&0&-\gamma\end{pmatrix}.
\end{equation}

\begin{proposition}[Characteristic polynomial]
\label{prop:char}
The characteristic polynomial of~\eqref{eq:J} is independent of helicity and equals
\begin{equation}
\label{eq:char}
p_q(\lambda)=(\lambda+\gamma)\bigl((\lambda+2H)^2+q^2\bigr)-A(\lambda+2H).
\end{equation}
In particular,
\begin{equation}
\label{eq:char0}
p_q(0)=\gamma(4H^2+q^2)-2HA.
\end{equation}
\end{proposition}

\begin{proof}
Direct expansion of $\det(\lambda I-J)$; the helicity terms cancel in the characteristic
polynomial.
Evaluating at $\lambda=0$ gives~\eqref{eq:char0}.
\end{proof}

\begin{theorem}[Zero-mode stationary threshold]
\label{thm:zero-mode}
Assume $H\neq0$.
Then
\begin{equation}
\label{eq:threshold}
p_0(0)=0\quad\Longleftrightarrow\quad A=2H\gamma.
\end{equation}
\end{theorem}

\begin{proof}
Set $q=0$ in~\eqref{eq:char0}.
\end{proof}

Define the stationary electromagnetic gain
\begin{equation}
\label{eq:GEM}
G_{\mathrm{EM}}(q)=\frac{A}{\gamma\bigl(2H+q^2/(2H)\bigr)}
\end{equation}
and the stationary critical relaxation rate
\begin{equation}
\label{eq:gammac}
\gamma_c(q)=\frac{2HA}{4H^2+q^2}.
\end{equation}

\begin{theorem}[Infrared-first stationary criticality]
\label{thm:IR-first}
Assume $A,H,\gamma>0$.
Then $G_{\mathrm{EM}}(q)$ is strictly decreasing in $q^2$, and
\begin{equation}
\gamma_c(q^2)\le\gamma_c(0)=\frac{A}{2H},
\end{equation}
with equality only at $q=0$.
Equivalently, among stationary zero-eigenvalue crossings, the $q=0$ mode reaches criticality
first as $\gamma$ decreases (or $A$ increases).
\end{theorem}

\begin{proof}
Differentiating~\eqref{eq:GEM} in $q^2$ yields a strictly negative quantity for positive
parameters.
The comparison for $\gamma_c$ is immediate from~\eqref{eq:gammac} since the denominator is
minimised at $q=0$.
\end{proof}

\begin{theorem}[Infrared unstable band]
\label{thm:band}
Assume $H,\gamma>0$.
Then
\begin{equation}
\label{eq:band}
p_q(0)<0
\quad\Longleftrightarrow\quad
q^2<\frac{2HA}{\gamma}-4H^2.
\end{equation}
In particular, if $A>2H\gamma$, the set of $q$ with $p_q(0)<0$ is a nonempty open infrared band
$0\le q<q_c$ with
\begin{equation}
\label{eq:qc}
q_c^2=\frac{2HA}{\gamma}-4H^2>0.
\end{equation}
Since $p_q(\lambda)\to+\infty$ as $\lambda\to+\infty$, the sign test $p_q(0)<0$ guarantees a
positive real root of $p_q$.
\end{theorem}

\begin{proof}
Rearrange~\eqref{eq:char0}.
The limit $p_q(\lambda)\to+\infty$ as $\lambda\to+\infty$ is immediate from the cubic leading
term, so a negative value at $\lambda=0$ forces a positive real root by the intermediate-value
theorem.
\end{proof}

\begin{remark}[Stationary versus Hopf]
\label{rem:hopf}
The zero-eigenvalue test does not by itself exclude an earlier Hopf crossing for arbitrary
parameter paths.
Theorem~\ref{thm:IR-first} is a stationary-branch statement: along the common experiment in
which $\gamma$ decreases through the zero-mode threshold, infrared modes go unstable first.
A global first-instability theorem for all paths requires the cubic Routh--Hurwitz criterion
on~\eqref{eq:char}; we do not claim that stronger result here.
\end{remark}

\section{Near-threshold growth}
\label{sec:growth}
Write $A=2H\gamma(1+\epsilon)$ with $0<\epsilon\ll1$.
Taylor expansion of $p_q(\lambda)=0$ about $(\lambda,q^2,\epsilon)=(0,0,0)$ yields the leading
dispersion
\begin{equation}
\label{eq:lambda}
\lambda(q)\simeq\frac{\gamma}{2H(2H+\gamma)}\bigl(4H^2\epsilon-q^2\bigr).
\end{equation}
Thus the leading unstable band is $q<2H\sqrt{\epsilon}$, and the maximum growth rate occurs at
$q=0$.

\begin{proposition}[Infrared onset, algebraic form]
\label{prop:onset}
Under the linearised threshold parametrisation $A=2H\gamma(1+\epsilon)$ and the frozen
denominator approximation of~\eqref{eq:lambda}, one has $\lambda(q)>0$ if and only if
$q^2<4H^2\epsilon$, with maximal growth at $q=0$.
\end{proposition}

\begin{proof}
Immediate from the sign of~\eqref{eq:lambda} for positive $\gamma,H$.
\end{proof}

This is an infrared onset rather than an ultraviolet catastrophe: the linear instability
selects the longest available physical wavelengths first.

\section{Nonlinear gauge-invariant completion}
\label{sec:nonlinear}
\begin{model}[Scalar--Maxwell kinetic completion]
\label{mod:I}
A minimal covariant completion uses a scalar response $h$ and a field-dependent Maxwell
kinetic function,
\begin{equation}
\label{eq:action}
S=\int\dd^4x\sqrt{-g}\left[
-\tfrac12(\nabla h)^2-V(h)-\tfrac14 I(h)^2F_{\mu\nu}F^{\mu\nu}
\right].
\end{equation}
A post-critical Landau potential
\begin{equation}
V(h)=-\tfrac12\mu^2 h^2+\tfrac{\lambda_h}{4}h^4,
\qquad\mu^2,\lambda_h>0,
\end{equation}
saturates at $|h_\star|=\mu/\sqrt{\lambda_h}$ in the absence of strong backreaction.
\end{model}

In conformal FLRW coordinates and Coulomb gauge, a transverse Fourier mode with canonical
variable $u_k=I A_k$ obeys
\begin{equation}
\label{eq:pump}
u_k''+\Bigl(k^2-\frac{I''}{I}\Bigr)u_k=0.
\end{equation}
Consequently a non-adiabatic homogeneous response is an electromagnetic pump.
A sufficient interval-wise condition for tachyonic growth is $I''/I>k^2$.
If $h$ oscillates about its minimum, the periodic part of $I''/I$ produces a Hill equation and
may generate parametric resonance.

\begin{remark}[What the pump does not prove]
\label{rem:pump}
Equation~\eqref{eq:pump} identifies a concrete channel from an infrared response instability to
finite-$k$ photon occupation.
It does not prove reheating.
Efficiency depends on the profile of $I(h)$, expansion, rescattering, strong-coupling
constraints, and backreaction.
The electromagnetic energy must be computed self-consistently with the $h$ equation and with
Einstein dynamics.
Those steps are modelling closures and open dynamical problems, not theorems of this paper.
\end{remark}

\section{Reciprocal conformal matching}
\label{sec:conformal}
Let a regular ``bridge'' metric $\hat g$ describe a neighbourhood of a spacelike crossover
$\Sigma$, and write
\begin{equation}
\label{eq:conf}
g^-_{\mu\nu}=\Omega_-^2\hat g_{\mu\nu},
\qquad
g^+_{\mu\nu}=\Omega_+^2\hat g_{\mu\nu}.
\end{equation}
Assume $\Omega_-\to\infty$ and $\Omega_+\to0$ as the crossover is approached from the outgoing
and incoming sides respectively, and impose the reciprocal ansatz
\begin{equation}
\label{eq:reciprocal}
\Omega_+\Omega_-=C_0,
\qquad C_0>0.
\end{equation}
This is the geometric skeleton shared with Penrose CCC~\cite{Penrose2006,Penrose2010,Tod2003}:
future conformal infinity of one aeon is identified with the Big Bang of the next.

\begin{assumption}[Four-dimensional Maxwell conformal transport]
\label{ass:maxwell}
In four spacetime dimensions, source-free Maxwell equations are conformally invariant when the
two-form $F_{\mu\nu}$ is held fixed~\cite{Wald1984}.
The covariant stress tensor and measured density then scale as
\begin{equation}
\label{eq:Tscale}
T_{\mu\nu}=\Omega^{-2}\hat T_{\mu\nu},
\qquad
\rho=\Omega^{-4}\hat\rho.
\end{equation}
\end{assumption}

\begin{theorem}[Cold outgoing / hot incoming]
\label{thm:hot-cold}
Under Assumption~\ref{ass:maxwell} and reciprocity~\eqref{eq:reciprocal}, finite nonzero
conformal radiation density $\hat\rho>0$ satisfies
\begin{equation}
\label{eq:hotcold}
\rho_-=\Omega_-^{-4}\hat\rho\to0,
\qquad
\rho_+=\Omega_+^{-4}\hat\rho\to\infty
\end{equation}
as $\Omega_-\to\infty$ and $\Omega_+\to0$ with $\Omega_+\Omega_-=C_0$.
Equivalently, with unit reciprocity $a_+a_-=1$ and fixed $\hat\rho\neq0$,
\begin{equation}
\label{eq:ratio8}
\frac{\rho_+}{\rho_-}=a_-^{8}.
\end{equation}
\end{theorem}

\begin{proof}
The scalings~\eqref{eq:Tscale} give $\rho_\pm=\Omega_\pm^{-4}\hat\rho$.
Reciprocity forces $\Omega_+=C_0/\Omega_-$, hence
$\rho_+/\rho_-=(\Omega_-/\Omega_+)^4=\Omega_-^{8}/C_0^{4}$.
In FLRW units with $a_\pm\propto\Omega_\pm$ and $C_0=1$, this is~\eqref{eq:ratio8}.
\end{proof}

\begin{remark}[Interpretation]
\label{rem:hotcold}
Theorem~\ref{thm:hot-cold} is a relation between two physical metrics and common conformal
data.
It must not be read as ordinary energy amplification inside a single metric.
On the outgoing side the same conformal radiation looks arbitrarily cold; on the incoming side
it looks arbitrarily hot.
That hot incoming branch is the big-bang-like content of the next aeon.
\end{remark}

\begin{theorem}[Reciprocal de~Sitter $\to$ radiation asymptotics]
\label{thm:radiation}
Assume the outgoing asymptotic~\eqref{eq:dS} and unit reciprocity $a_+(\eta)=1/a_-(-\eta)$ for
$\eta>0$ small.
Then
\begin{equation}
\label{eq:aPlus}
a_+(\eta)=H_-\eta+o(\eta).
\end{equation}
Writing $\dd t_+=a_+\dd\eta$ and $a_+=C\eta$ with $C=H_-$, one obtains
\begin{equation}
\label{eq:radiation}
a_+^{2}=2C\,t_+,
\end{equation}
hence the radiation-era FLRW exponent $a_+\propto t_+^{1/2}$.
\end{theorem}

\begin{proof}
Substitute~\eqref{eq:dS} into the reciprocal ansatz.
With $a=C\eta$ and $t=\int a\,\dd\eta=C\eta^{2}/2$ one has $a^{2}=2Ct$ by direct algebra.
\end{proof}

\begin{proposition}[Reciprocal physical momenta]
\label{prop:mom}
If $a_+a_-=1$ and $a_-\neq0$, then $k/a_+=k\,a_-$ for any comoving wavenumber $k$.
Soft outgoing physical momenta are therefore hard on the incoming side.
\end{proposition}

\begin{proof}
Immediate from $1/a_+=a_-$.
\end{proof}

\begin{remark}[Normalization caveat]
\label{rem:norm}
Reciprocity alone fixes an asymptotic form only up to normalization.
If one additionally assumes a flat radiation-dominated Friedmann equation and writes
$a_+=C\eta$, then $\hat\rho_\gamma=3C^{2}/(8\pi G)$.
Setting $C=H_-$ by unit reciprocity gives $\hat\rho_\gamma=3H_-^{2}/(8\pi G)$.
That equality is a consistency condition involving conformal gauge and Einstein dynamics, not
an independent prediction of Maxwell conformal invariance alone.
\end{remark}

\section{Active versus passive conformal cyclic cosmology}
\label{sec:active}
Penrose's CCC proposes an infinite succession of aeons in which the conformal infinity of one
aeon is identified with the Big Bang of the next~\cite{Penrose2006,Penrose2010,Penrose2013CCC}.
The geometric matching---a spacelike crossover, conformal factors diverging/vanishing, and
transport of zero-rest-mass field data---is the shared spine of CCC and HCCC.

The dynamical reading differs.
\begin{center}
\begin{tabular}{@{}>{\bfseries}p{0.22\linewidth}p{0.35\linewidth}p{0.35\linewidth}@{}}
\toprule
& Penrose CCC & HCCC (this paper)\\
\midrule
Geometry & Conformal aeon bridge & Same reciprocal conformal bridge\\
Driver & Passive: massless conformal era after matter evaporates & Active: IR EM Hysterical phase transition\\
Radiation seed & Inherited conformal data / Hawking evaporation products & Photon pump from supercritical response~\eqref{eq:pump}\\
Big Bang of next aeon & Conformally reinterpreted future infinity & Hot incoming branch of Theorem~\ref{thm:hot-cold}\\
Status claimed here & Literature geometry & Geometry + IR instability theorems; pump/junction open\\
\bottomrule
\end{tabular}
\end{center}

\begin{definition}[Active big-bang-like crossover]
\label{def:active-bang}
A crossover $\Sigma$ is \emph{active big-bang-like} when:
\begin{enumerate}
\item an outgoing aeon approaches a de~Sitter-like conformal infinity;
\item an instability produces nonzero conformal radiation data $\hat\rho$ on a neighbourhood of
$\Sigma$;
\item reciprocal conformal matching maps those data to an incoming radiation-dominated
asymptotics with $\rho_+\to\infty$ as in Theorem~\ref{thm:hot-cold}.
\end{enumerate}
\end{definition}

Under Modelling Assumption~\ref{mod:I}, Assumption~\ref{ass:maxwell}, and the reciprocal ansatz,
the infrared electromagnetic mechanism of Sections~\ref{sec:linear}--\ref{sec:conformal}
supplies a concrete candidate for Definition~\ref{def:active-bang}.
The phase transition is the physical event; the conformal matching is the geometric
translation of that event into the language of the next aeon's Big Bang.

\begin{remark}[What is and is not identified with ``our'' Big Bang]
\label{rem:ourBB}
HCCC claims that \emph{some} aeon-to-aeon crossover of this type is big-bang-like for the
incoming side.
Whether the Big Bang of our observed aeon was produced by a previous Hysterical electromagnetic
crossover is a physical hypothesis, not a theorem.
Likewise, observational CCC proposals (concentric low-variance circles, Hawking points) are not
evaluated here~\cite{Penrose2010,Penrose2013CCC}.
\end{remark}

\section{Logical status of the results}
\label{sec:status}
\begin{center}
\small
\begin{tabular}{@{}p{0.48\linewidth}p{0.44\linewidth}@{}}
\toprule
\textbf{Quantity / claim} & \textbf{Status}\\
\midrule
Characteristic polynomial~\eqref{eq:char} from $\det(\lambda I-J)$ & Prop.~\ref{prop:char} (Lean: \texttt{prop\_char}, helicity $s^2=1$).\\
Zero-mode threshold $A=2H\gamma$ & Theorem~\ref{thm:zero-mode} (Lean).\\
IR-first stationary criticality (incl.\ strict) & Theorem~\ref{thm:IR-first} (Lean: \texttt{thm\_IR\_first\_full}).\\
Infrared unstable band / sign test / positive root & Theorem~\ref{thm:band} (Lean: \texttt{thm\_band\_full}, IVT).\\
Near-threshold dispersion sign & Prop.~\ref{prop:onset} (Lean algebra of the frozen-denominator form).\\
Scalar--Maxwell action / $I(h)$ pump & Modelling Assumption~\ref{mod:I}; not Lean.\\
Photon production efficiency / reheating & Open dynamical problem.\\
4D Maxwell conformal invariance & Literature hypothesis~\cite{Wald1984}.\\
Cold/hot density dichotomy; $\rho_+/\rho_-=a_-^{8}$ & Theorem~\ref{thm:hot-cold} (Lean: unit and $C_0$ forms).\\
Reciprocal dS$\to$radiation (exact leading / linear ansatz) & Theorem~\ref{thm:radiation} (Lean: \texttt{thm\_radiation\_full}; little-$o$ remainder not re-proved).\\
Reciprocal momenta & Prop.~\ref{prop:mom} (Lean).\\
Reciprocal conformal ansatz $\Omega_+\Omega_-=C_0$ & Modelling / CCC geometry input.\\
Saturation value $\hat\rho_{\mathrm{EM}}(\Sigma)$ & Not derived.\\
Einstein--Maxwell--response junction regularity & Open.\\
Entropy production / thermalisation & Open.\\
Observational CCC signatures & Not claimed.\\
``Nature realises HCCC'' & Not claimed.\\
\bottomrule
\end{tabular}
\end{center}

\section{Conclusion}
\label{sec:conclusion}
The late-universe gravitational package of this series leaves an electromagnetic window:
finite-wavenumber response can become critical even though homogeneous net charge cannot be
created from neutrality.
In an asymptotically de~Sitter background, a minimal helical electromagnetic--response model
has an infrared-first stationary instability and a finite unstable band.
A scalar-dependent Maxwell kinetic function then offers a gauge-invariant route from that
infrared onset to finite-$k$ photons.
Reciprocal conformal matching---the geometric core shared with Penrose's conformal cyclic
cosmology---maps outgoing de~Sitter asymptotics to an incoming radiation-era law and converts
regular conformal radiation data into a hot big-bang-like branch for the next aeon.

The resulting Hysterical Conformal Cyclic Cosmology is therefore CCC-like in conformal
geometry and distinct in dynamics: the crossover is an active phase transition, not a passive
conformal identification after the last massive degrees of freedom fade.
The theorems close the algebraic spine of that claim.
The scientific burden that remains is dynamical: derive the saturated conformal radiation
amplitude, solve the coupled response--Maxwell--Einstein system through the junction, and
demonstrate approximate thermalisation.
Until those steps are closed, HCCC is a mathematically structured mechanism for an active
aeon bridge, not a completed cyclic cosmology of Nature.

\bibliographystyle{plain}
\bibliography{paper_refs}
\end{document}
